\chapter{多工况仿真结果分析与不利工况识别}

\section{工况变化规律分析}

在前文模型校核、工况矩阵设计基础之上，本文就平坦地形、起伏地形、山脊地形、谷地地形和阶跃地形这五种基础工况展开对比分析。不同工况时飞行器着陆末段的离地高度、速度、姿态角和横向稳定性存在差异，证明地形条件对固定翼飞行器着陆控制性能有直接影响。

从离地高度的变化规律来看，在平坦地形工况下，飞行器的AGL曲线变化最为平缓，在整个着陆过程中沿着预定的下滑轨迹持续下降，末段控制过程比较稳定。在起伏地形工况下，因为地面高程一直在变化，飞行器相对离地高度呈现出更为明显的非线性变化，使得系统在末段需要进行更为频繁的高度修正。在山脊地形工况下，飞行器接近跑道区域之前地面高程快速抬升，致使相对离地高度在短时间内明显减小，这对拉平时机提出了更高的要求。谷地地形工况表现为前段高度裕度较大而后段相对离地高度变化加快的特点，增加了后段控制节奏的复杂性。在阶跃地形工况中，地面高程存在突变，让飞行器AGL变化出现明显的压缩特征，对末段接地判断的影响最大。

就速度保持规律而言，在各个工况下飞行器末段的速度整体上都被控制在了安全范围之内，表明系统大体上能够达成速度约束这项任务。不过不同地形条件下速度调节的难度存在差异。在平坦地形和谷地地形工况中，速度变化相对比较平稳，末段速度保持效果比较理想；在起伏地形和山脊地形工况中，因为要同时兼顾高度修正和姿态调整，速度波动有一定程度的增加；在阶跃地形工况中，虽然速度没有显著超出界限，但由于近地阶段高度变化更为突然，速度控制与姿态控制之间协调的压力明显更大。

依据姿态变化规律来分析，在平坦地形工况中，俯仰角和滚转角的变化幅度最小，意味着系统在标准着陆条件下能够维持较好的纵横向稳定性。在起伏地形与山脊地形工况中，飞行器为了持续对下滑状态和跑道对准进行修正，俯仰角和滚转角的波动更加明显。在谷地地形工况中，俯仰角整体比较平顺，但滚转角呈现出一定的偏转特征，表明其横向稳定性受到了地形的影响。在阶跃地形工况中，姿态修正最为频繁，体现出系统需要在较短时间内同时兼顾高度、速度和方向的控制。地形越是复杂，着陆末段状态变化就越敏感，系统对于操纵辅助和风险提示的需求也就越强。

\section{不利工况与较优工况识别}

按照之前所建立的评价指标与判据，能够对五类工况的优劣予以识别。从末段AGL、末段速度、俯仰角、滚转角、控制过程平顺性这几个方面进行综合考量，平坦地形工况的整体表现最为出色，能够作为较优工况。在该工况下，飞行器沿着预定的轨迹稳定下降，末段离地高度约为4.60 ft，末段速度约为71.80 kts，俯仰角和滚转角都处于相对稳定的水平，表明系统在标准环境中具备较好的收敛性与稳定性。

谷地地形工况的仿真总时长比较长，不过末段速度保持的效果不错，末段离地高度和平坦地形相近，俯仰角的变化也比较小，总体性能比较理想，可当作次优工况。该工况的主要特点是前段控制裕度大，系统能够较好地进行轨迹修正，末段状态总体比较平顺。

起伏地形工况和山脊地形工况能够被归结为偏不利的工况类型。在起伏地形工况中，俯仰角和滚转角的变化幅度相对较大，意味着系统需要持续对处于连续变化状态的地面高程作出响应。在山脊地形工况下，鉴于地形在接近区域呈现快速抬升的趋势，飞行器相对离地高度的减小速度更为迅速，倘若控制修正稍微出现滞后，便极易对拉平质量产生影响。这两种工况虽然并未导致明显的失稳现象，但已经暴露出系统在复杂环境下的薄弱环节。

在各种工况中，阶跃地形工况可被判定为最为不利的工况。其最为突出的特征在于末段AGL仅为1.20 ft，这一数值明显低于其他工况，说明地形发生突变后近地阶段的控制调整空间被极大幅度压缩。该工况末段滚转角的绝对值相对较大，表明横向姿态保持的难度也随之增大。虽然速度没有出现严重越界，但高度突变对操纵者判断拉平时机和接地窗口造成了比较明显的干扰。因此，阶跃地形工况是需要重点关注的对象。

\section{典型工况机理解释}

平坦地形工况能够呈现比较优异的工况表现，其主要缘由在于地面高程基本维持恒定状态，飞行器海拔高度与相对离地高度之间的关系清晰简单，系统能够依据预先设定的下滑角、目标速度和阶段切换规则实现稳定运行，操纵者对末段接地时机的判断也更易于与系统提示达成一致，整体控制最容易趋向收敛。

山脊地形工况所存在的问题主要是因为接近跑道之前地面高程迅速上升。飞行器海拔高度虽然依旧平稳下降，但相对离地高度会在较短时间内急剧减小，使得系统必须更早地完成下沉率抑制和姿态修正。修正速度偏慢容易导致拉平不足，修正过度又可能致使姿态波动增加。山脊地形工况体现出系统对地形快速变化比较敏感。

阶跃地形工况被认为是最不利的，其根本缘由是地面高程的不连续变化使得飞行器相对离地高度突然降低。对操纵者来说这种情况会压缩反应时间，对控制系统而言则会加大阶段识别、边界判断和接地时机控制的难度。这种工况并非只是纵向控制困难，而是高度判断、姿态控制、接地窗口识别同时变得困难。因此在复杂地形特别是阶跃型地形条件下，仅仅依靠基础控制律难以充分确保着陆品质，还需要借助更清晰的状态提示、风险告警和约束保护来提高系统鲁棒性。

综合来看，通过多工况分析能够发现，平坦地形可作为系统的基准参考，谷地地形整体呈现较好状况，起伏、山脊和阶跃地形更能体现系统在复杂环境下的风险特性。其中阶跃地形属于最为典型的不利工况，应作为后续人机交互强化设计和系统优化时重点分析的对象。
